\documentclass[10pt,a4paper]{amsart}
\usepackage{amsmath,amssymb,amsthm,graphicx}
\usepackage{verbatim}
\newenvironment{reference}{\comment}{\endcomment}
%\newenvironment{reference}{}{}
\newenvironment{slogan}{\comment}{\endcomment}
\newenvironment{history}{\comment}{\endcomment}

% For commutative diagrams we use Xy-pic
\usepackage[all]{xy}

% We use 2cell for 2-commutative diagrams.
\xyoption{2cell}
% Theorem environments.
%
\theoremstyle{plain}
\newtheorem{theorem}[subsection]{Theorem}
\newtheorem{proposition}[subsection]{Proposition}
\newtheorem{lemma}[subsection]{Lemma}

\theoremstyle{definition}
\newtheorem{definition}[subsection]{Definition}
\newtheorem{example}[subsection]{Example}
\newtheorem{exercise}[subsection]{Exercise}
\newtheorem{situation}[subsection]{Situation}

\theoremstyle{remark}
\newtheorem{remark}[subsection]{Remark}
\newtheorem{remarks}[subsection]{Remarks}

\numberwithin{equation}{subsection}

% Macros
%
\def\lim{\mathop{\mathrm{lim}}\nolimits}
\def\colim{\mathop{\mathrm{colim}}\nolimits}
\def\Spec{\mathop{\mathrm{Spec}}}
\def\Hom{\mathop{\mathrm{Hom}}\nolimits}
\def\Ext{\mathop{\mathrm{Ext}}\nolimits}
\def\SheafHom{\mathop{\mathcal{H}\!\mathit{om}}\nolimits}
\def\SheafExt{\mathop{\mathcal{E}\!\mathit{xt}}\nolimits}
\def\Sch{\mathit{Sch}}
\def\Mor{\mathop{\mathrm{Mor}}\nolimits}
\def\Ob{\mathop{\mathrm{Ob}}\nolimits}
\def\Sh{\mathop{\mathit{Sh}}\nolimits}
\def\NL{\mathop{N\!L}\nolimits}
\def\CH{\mathop{\mathrm{CH}}\nolimits}
\def\proetale{{pro\text{-}\acute{e}tale}}
\def\etale{{\acute{e}tale}}
\def\QCoh{\mathit{QCoh}}
\def\Ker{\mathop{\mathrm{Ker}}}
\def\Im{\mathop{\mathrm{Im}}}
\def\Coker{\mathop{\mathrm{Coker}}}
\def\Coim{\mathop{\mathrm{Coim}}}

% Boxtimes
%
\DeclareMathSymbol{\boxtimes}{\mathbin}{AMSa}{"02}

%
% Macros for moduli stacks/spaces
%
\def\QCohstack{\mathcal{QC}\!\mathit{oh}}
\def\Cohstack{\mathcal{C}\!\mathit{oh}}
\def\Spacesstack{\mathcal{S}\!\mathit{paces}}
\def\Quotfunctor{\mathrm{Quot}}
\def\Hilbfunctor{\mathrm{Hilb}}
\def\Curvesstack{\mathcal{C}\!\mathit{urves}}
\def\Polarizedstack{\mathcal{P}\!\mathit{olarized}}
\def\Complexesstack{\mathcal{C}\!\mathit{omplexes}}
% \Pic is the operator that assigns to X its picard group, usage \Pic(X)
% \Picardstack_{X/B} denotes the Picard stack of X over B
% \Picardfunctor_{X/B} denotes the Picard functor of X over B
\def\Pic{\mathop{\mathrm{Pic}}\nolimits}
\def\Picardstack{\mathcal{P}\!\mathit{ic}}
\def\Picardfunctor{\mathrm{Pic}}
\def\Deformationcategory{\mathcal{D}\!\mathit{ef}}
\usepackage{cleveref}
\crefname{originallemma}{original lemma}{original lemmas}
\Crefname{originallemma}{Original lemma}{Original lemmas}
\makeatletter
\newlabel{etale-cohomology-lemma-constructible-quasi-compact-quasi-separated}{{095E}{source}{}{sourcepointer-1}{}}
\newlabel{etale-cohomology-lemma-constructible-quasi-compact-quasi-separated@cref}{{[originallemma][0][]095E}{[1][0][]source}{}{}{}}
\makeatother
\crefname{originalproposition}{original proposition}{original propositions}
\Crefname{originalproposition}{Original proposition}{Original propositions}
\makeatletter
\newlabel{etale-cohomology-proposition-constructible-over-noetherian}{{09BH}{source}{}{sourcepointer-2}{}}
\newlabel{etale-cohomology-proposition-constructible-over-noetherian@cref}{{[originalproposition][0][]09BH}{[1][0][]source}{}{}{}}
\makeatother
\crefname{originalreference}{original reference}{original references}
\Crefname{originalreference}{Original reference}{Original references}
\makeatletter
\newlabel{lemma-compare-constructible}{{09AK}{source}{}{sourcepointer-3}{}}
\newlabel{lemma-compare-constructible@cref}{{[originalreference][0][]09AK}{[1][0][]source}{}{}{}}
\makeatother
\crefname{originallemma}{original lemma}{original lemmas}
\Crefname{originallemma}{Original lemma}{Original lemmas}
\makeatletter
\newlabel{sites-modules-lemma-kernel-finite-locally-constant}{{093U}{source}{}{sourcepointer-4}{}}
\newlabel{sites-modules-lemma-kernel-finite-locally-constant@cref}{{[originallemma][0][]093U}{[1][0][]source}{}{}{}}
\makeatother
\title{A constructible sheaf on the pro-etale site of a Noetherian scheme is lisse on the pieces of a finite partition}
\author{The Stacks Project Authors}
\date{stacks-project:proetale.tex:local-snapshot-20261012; GNU FDL 1.2 or later}
\begin{document}
\maketitle
\noindent\textit{Editorial note.} All delivered mathematical text is the author's own, copied from the pinned Stacks Project chapter file proetale.tex; no statement or proof step was rewritten, reordered or removed.\par\medskip
\noindent\textit{Editorial note.} The counted proof is the author's complete argument: the construction of the associated graded sheaf over the Noetherian ring of quotients of powers of the ideal together with its surjectivity from degree zero, the partition making that graded sheaf locally constant, and the induction along the short exact sequences propagating local constancy back to the original sheaf.\par\medskip
\noindent\textit{Editorial note.} The label lemma-Noetherian-constructible names the counted statement, and the proof environment immediately following it in the source is delivered verbatim as its proof.\par\medskip
\noindent\textit{Editorial note.} The counted statement is carried by the author's own lemma environment in the chapter source.\par\medskip
\noindent\textit{Editorial note.} Pre-open gate: python operations/agent-authoring/tools/precard\_gate.py --source proetale.tex --label lemma-Noetherian-constructible --start 5731 --end 5780 --local-id 3585 returned RESULT=OK (L3585 free, occupancy FREE, public interval clear). Candidate taken from sources/stacks-proetale/candidates-halfA.json.\par\medskip
\noindent\textit{Editorial note.} The statement and proof are copied from the pinned chapter file; the cited inputs are kept as external pointers with their canonical Stacks tags.\par\medskip
\noindent\textit{Allowed standard prerequisites (editorial).}\begin{itemize}
\item Comparison of constructible sheaves on the pro-etale and etale sites. Same chapter, the lemma identifying constructible sheaves on the pro-etale site with the pullbacks of constructible etale sheaves. Supplies the existence of the constructible etale sheaves G\_n whose pullbacks are the associated graded pieces I\textasciicircum{}n F / I\textasciicircum{}\{n+1\} F.
\item Constructibility of a sheaf of modules over a Noetherian graded ring generated by its degree zero part. Etale Cohomology, the proposition that a sheaf of R-modules generated by the pullback of a constructible sheaf is constructible over a Noetherian scheme. Supplies the constructibility of the direct sum G of the G\_n as a sheaf of R-modules.
\item Existence of a partition on which a constructible sheaf of finite type is locally constant. Etale Cohomology, the lemma producing a finite partition into locally closed subschemes on which a constructible sheaf of finite type is locally constant. Supplies the partition by locally closed subschemes on which G becomes locally constant.
\item Local constancy of the kernel of a surjection of locally constant sheaves of finite type. Modules on Sites, the lemma that the kernel of a surjection between locally constant sheaves of finite type is again locally constant of finite type. Supplies the induction step showing that the successive quotients I\textasciicircum{}n F / I\textasciicircum{}\{n+1\} F are locally constant, and hence that F is lisse.
\end{itemize}\medskip
\ifdefined\hypertarget\hypertarget{sourcepointer-1}{}\fi
\noindent\textit{Source reference.} Original lemma 095E: sources/stacks-chow/tags.txt, source lines 11582--11582. Prerequisite pointer invoked at proof line(s) 5767; it names a standard fact of a sibling chapter that the counted argument uses, and every step of the counted proof built on it is retained verbatim in the delivered body.\par\medskip
\ifdefined\hypertarget\hypertarget{sourcepointer-2}{}\fi
\noindent\textit{Source reference.} Original proposition 09BH: sources/stacks-chow/tags.txt, source lines 11795--11795. Prerequisite pointer invoked at proof line(s) 5763; it names a standard fact of a sibling chapter that the counted argument uses, and every step of the counted proof built on it is retained verbatim in the delivered body.\par\medskip
\ifdefined\hypertarget\hypertarget{sourcepointer-3}{}\fi
\noindent\textit{Source reference.} Original reference 09AK: sources/stacks-proetale/pinned-original-source.tex, source lines 5456--5456. Prerequisite pointer invoked at proof line(s) 5748; it names a standard fact of the same chapter that the counted argument uses, and every step of the counted proof built on it is retained verbatim in the delivered body.\par\medskip
\ifdefined\hypertarget\hypertarget{sourcepointer-4}{}\fi
\noindent\textit{Source reference.} Original lemma 093U: sources/stacks-chow/tags.txt, source lines 11527--11527. Prerequisite pointer invoked at proof line(s) 5778; it names a standard fact of a sibling chapter that the counted argument uses, and every step of the counted proof built on it is retained verbatim in the delivered body.\par\medskip
\begin{lemma}
\label{lemma-Noetherian-constructible}
Let $X$ be a Noetherian scheme. Let $\Lambda$ be a Noetherian ring and
let $I \subset \Lambda$ be an ideal. Let $\mathcal{F}$ be a
constructible $\Lambda$-sheaf on $X_\proetale$.
Then there exists a finite partition $X = \coprod X_i$ by
locally closed subschemes such that the restriction $\mathcal{F}|_{X_i}$
is lisse.
\end{lemma}

\begin{proof}
Let $R = \bigoplus I^n/I^{n + 1}$. Observe that $R$ is a Noetherian ring.
Since each of the sheaves
$\mathcal{F}/I^n\mathcal{F}$ is a constructible sheaf of
$\Lambda/I^n\Lambda$-modules also $I^n\mathcal{F}/I^{n + 1}\mathcal{F}$
is a constructible sheaf of $\Lambda/I$-modules and hence the pullback
of a constructible sheaf $\mathcal{G}_n$ on $X_\etale$ by
Lemma \ref{lemma-compare-constructible}.
Set $\mathcal{G} = \bigoplus \mathcal{G}_n$. This is a sheaf
of $R$-modules on $X_\etale$ and the map
$$
\mathcal{G}_0 \otimes_{\Lambda/I} \underline{R}
\longrightarrow
\mathcal{G}
$$
is surjective because the maps
$$
\mathcal{F}/I\mathcal{F} \otimes \underline{I^n/I^{n + 1}} \to
I^n\mathcal{F}/I^{n + 1}\mathcal{F}
$$
are surjective. Hence $\mathcal{G}$ is a constructible sheaf of
$R$-modules by \'Etale Cohomology, Proposition
\ref{etale-cohomology-proposition-constructible-over-noetherian}.
Choose a partition $X = \coprod X_i$ such that
$\mathcal{G}|_{X_i}$ is a locally constant sheaf of $R$-modules
of finite type (\'Etale Cohomology, Lemma
\ref{etale-cohomology-lemma-constructible-quasi-compact-quasi-separated}).
We claim this is a partition as in the lemma.
Namely, replacing $X$ by $X_i$ we may assume $\mathcal{G}$ is locally
constant. It follows that each of the sheaves
 $I^n\mathcal{F}/I^{n + 1}\mathcal{F}$
is locally constant. Using the short exact sequences
$$
0 \to I^n\mathcal{F}/I^{n + 1}\mathcal{F} \to
\mathcal{F}/I^{n + 1}\mathcal{F} \to \mathcal{F}/I^n\mathcal{F} \to 0
$$
induction and Modules on Sites, Lemma
\ref{sites-modules-lemma-kernel-finite-locally-constant}
the lemma follows.
\end{proof}

\end{document}
